\documentclass[journal]{IEEEtran}
\usepackage[margin=1in]{geometry}
\pdfinfoomitdate=1
\pdftrailerid{}
\pdfsuppressptexinfo=-1
\usepackage{amsmath,amssymb,amsthm,bm,booktabs}
\usepackage{xr-hyper}
\usepackage[hidelinks]{hyperref}
\externaldocument{main}[main.pdf]
\newcommand{\E}{\mathbb E}
\newcommand{\tr}{\operatorname{tr}}
\newcommand{\Pp}{\mathbb P}
\title{Supplementary Material: Learning-Aware Lyapunov Resource Control for Correlated Multi-Agent Markov Learning}
\author{Jialin Zhuang and Vincent K. N. Lau}
\begin{document}
\maketitle
\renewcommand{\theequation}{S\arabic{equation}}

\section{Proof of the Same-Stream Actual-Risk Enclosure}
Consider the stationary common factor in the main paper, with independent private observation noise and a deterministic probe $d=(q_d,b_d,n_d)$ ending at $T=n_db_d$.
The probe sigma-field $\mathcal H_d$ contains its observations; spatial contrasts contain no additional information about the common factor.
Under true instance $j$, Gaussian filtering gives
\begin{align}
C_T\mid\mathcal H_d&\sim\mathcal N(m_j,P_j),\nonumber\\
m_j&\sim\mathcal N(0,s_j^2),\quad s_j^2=\theta_j-P_j.
\end{align}
The posterior variance $P_j$ is deterministic because the probe schedule is fixed before data.
Starting from variance $\theta_j$, the observation update is $P\leftarrow P/(1+q_dP)$; between observations it is $P\leftarrow\lambda^{2b_d}P+(1-\lambda^{2b_d})\theta_j$.
The last update has no subsequent prediction.

Fix either downstream program $a$, including its budget, initialization, sample spacing and deterministic weights.
All learning observations occur strictly after $T$.
With $r_a=\lambda^{b_a}$ and the notation of main Eq.~\eqref{eq:linear-program}, its unconditional stationary risk is
\begin{align}
J_{j,a}
&=d_a^2+\theta_j\sum_{t,u=1}^{N_a}\omega_{a,t}\omega_{a,u}r_a^{|t-u|}\nonumber\\
&\quad+q_a^{-1}\sum_{t=1}^{N_a}\omega_{a,t}^2.
\label{eq:supp-exact-risk}
\end{align}
Conditional means of future common states are $\lambda^{b_at}m_j$.
Their conditional covariance is their stationary covariance minus $s_j^2\lambda^{b_a(t+u)}$.
Private learning noise is independent of the probe.
Therefore
\begin{equation}
\E_j[Z_a^2\mid\mathcal H_d]
=J_{j,a}+2d_ak_am_j+k_a^2(m_j^2-s_j^2).
\label{eq:supp-conditional-risk}
\end{equation}
This identity includes the posterior mean, which is unbounded across probe histories; no uniform conditional Gaussian bound is invoked.

Let $E_j$ be the event that the decision selects the wrong frozen program, and let $p_j=\Pr_j(E_j)\le\delta$.
Subtracting Eq.~\eqref{eq:supp-conditional-risk} for the correct and wrong programs gives
\begin{equation}
\E_jZ_{\widehat a}^2=J_{j,a_c}+p_j\Delta J_j+\E_j[\mathbf1_{E_j}U_j],
\end{equation}
where
\begin{align}
U_j&=2(d_{a_w}k_{a_w}-d_{a_c}k_{a_c})m_j\nonumber\\
&\quad+(k_{a_w}^2-k_{a_c}^2)(m_j^2-s_j^2).
\end{align}
The centered Gaussian first and second moments are orthogonal, with variances $s_j^2$ and $2s_j^4$.
Consequently $\E_jU_j=0$ and $\E_jU_j^2=\Gamma_j^2$ as defined in the main paper.
Cauchy--Schwarz gives $|\E_j[\mathbf1_{E_j}U_j]|\le\sqrt{p_j}\Gamma_j\le\sqrt\delta\Gamma_j$.
Bounding the signed term $p_j\Delta J_j$ and using nonnegative squared risk proves main Theorem~\ref{thm:same-stream}.

\section{Exact Weights of the Implemented Learners}
Let $M_a=\min\{\lfloor B_m'/(h+q_a)\rfloor,\lfloor B_e'/b_a\rfloor\}$ be the number of scheduled observations after subtracting probe costs.
All $M_a$ observations are paid.
The FIFO withholding convention leaves $N_a=(M_a-D)_+$ usable observations, not a free replacement budget.
For $N_a=0$ the reported parameter is one, so $d_a=1$ and the weight vector is empty.

\subsection{Running mean}
The Layer-A output is $(1+\sum_{t=1}^{N_a}\bar X_t)/(N_a+1)$.
Thus $d_a=1/(N_a+1)$ and every $\omega_{a,t}=1/(N_a+1)$.
Main Eq.~\eqref{eq:linear-program} then gives
\begin{equation}
k_a=\frac{r_a(1-r_a^{N_a})}{(N_a+1)(1-r_a)}.
\end{equation}

\subsection{Averaged FIFO affine learner}
In Layer B, an increment $\eta(\bar X_t-\mu w_{t-1})$ is computed before enqueueing and is applied after $D$ opportunities.
The original implementation fixes $\mu=1/2$ and $\eta=0.1/[\mu(D+1)]$, starts from parameter one with an empty queue, and averages the applied parameters.
Indexing these applied parameters as $z_n$ gives
\begin{equation}
z_n=z_{n-1}-\eta\mu z_{n-D-1}+\eta\bar X_n,
\quad z_{-D}=\cdots=z_0=1.
\end{equation}
Define $h_0=1$, $h_t=0$ for $t<0$, and $h_t=h_{t-1}-\eta\mu h_{t-D-1}$ for $t\ge1$.
For $N=N_a>0$, the output $N^{-1}\sum_{n=1}^Nz_n$ has exact coefficients
\begin{equation}
\omega_{a,t}=\frac\eta N\sum_{v=0}^{N-t}h_v,
\quad d_a=1-\mu\sum_{t=1}^N\omega_{a,t}.
\label{eq:supp-fifo-weights}
\end{equation}
To verify the deterministic term, constant input $\bar X_n=\mu$ keeps the initial history at one; linearity then yields the identity for $d_a$.
Eq.~\eqref{eq:supp-fifo-weights} represents the average of applied iterates; all scheduled observations remain charged.

\section{Shrinking-Risk and Budget-Shift Interface}
Suppose the finite action catalogue has $|\lambda|\le1-\gamma<1$, bounded $D$, and strictly stable affine recursions.
For the registered step rule this stability follows directly from the characteristic polynomial $z^{D+1}-z^D+t$, $t=\eta\mu$.
At small $t>0$, the root at one moves inside the unit disk and the other roots remain inside.
Writing a first unit-circle crossing as $z=e^{i\omega}$ yields $t=2\sin(\omega/2)$ and the smallest positive crossing $\omega=\pi/(2D+1)$.
Thus $0<t<2\sin[\pi/(4D+2)]$ is stable; the registered $t=0.1/(D+1)$ lies strictly inside this interval.
The impulse response consequently decays geometrically, and $\sum_{v\ge0}h_v=1/(\eta\mu)$.

The interior weights in Eq.~\eqref{eq:supp-fifo-weights} are $1/(\mu N)$ up to a summable end correction: their difference from these constant weights has $\ell_1$ norm $O(N^{-1})$ and maximum absolute value $O(N^{-1})$.
The homogeneous transient is summable, so its average $d_a$ is $O(N^{-1})$.
The AR(1) covariance kernel has uniformly bounded absolute row sums.
Substitution into Eq.~\eqref{eq:supp-exact-risk} therefore gives
\begin{equation}
J_{j,a}=\frac1{\mu^2N_a}\left[\frac1{q_a}+\theta_j\frac{1+r_a}{1-r_a}\right]+O(N_a^{-2}).
\label{eq:supp-averaged-rate}
\end{equation}
The running mean has the same expansion with $\mu=1$.
These constants are uniform over the finite catalogue and bounded instance class.
On a budget ray, $N_a=s/c_a+O(v+D+1)$ with $c_a=\max\{(h+q_a)/\beta_m,b_a/\beta_e\}$ and fixed probe charge $v$.
Hence $J_{j,a}=K_{j,a}/s+O((v+D+1)/s^2)$, with
\begin{equation}
K_{j,a}=\frac{c_a}{\mu^2}\left[\frac1{q_a}+\theta_j\frac{1+r_a}{1-r_a}\right].
\end{equation}
Subtracting a fixed probe charge changes $N_a$ by at most $O(v+D+1)$, giving the claimed budget-shift remainder.
Moreover $k_a=O(N_a^{-1})$ because the weights are uniformly $O(N_a^{-1})$ and $\sum_{t\ge1}|r_a|^t$ is bounded.
Since $s_j^2\le\theta_j$, main Eq.~\eqref{eq:selection-carryover} yields
\begin{equation}
\sqrt\delta\Gamma_j=O(\sqrt\delta/s^2).
\end{equation}
For a pair of frozen actions satisfying the main phase theorem's leading-gap and wrong-decision assumptions, adding this remainder to admission changes only its second-order class constants.
The same necessary-information and sufficient-scale argument then applies to the corrected actual-risk rule.
This is an averaged, additive Gaussian affine result; it does not remove the nonzero noise floor of a fixed-step unaveraged learner or the conditional-history requirements of the general Markov theorem.

\section{Complete Existing-Evidence Reanalysis}
The calculation uses all registered EXP-016B budget points, both regimes and all eight policies: 14,336 policy cells, each with its original 192-seed empirical mean.
Layer A contains 104 scenario families and Layer B its original 56, including eight search-censored families in each layer.
For fixed policies Eq.~\eqref{eq:supp-exact-risk} is exact.
For probe policies main Eq.~\eqref{eq:same-stream-risk} supplies population-risk bounds.
When information-only and learning-aware policies execute the same frozen program, their paired difference is identically zero, not the difference of two independent interval bounds.
Otherwise subtract their one-sided bounds.

\begin{table}[t]
\centering
\caption{Population enclosures for mean information-only minus learning-aware risk, rounded outward.
All means use the listed complete stratum.}
\begin{tabular}{llrrr}
\toprule
Layer & Stratum & Cells & Lower & Upper\\
\midrule
A & All & 1184 & .0506 & .0706\\
A & Intermediate & 576 & .0770 & .1180\\
A & Original practical & 446 & .1607 & .1952\\
A & Search-censored & 32 & .4879 & .4880\\
B & All & 608 & .0079 & .0223\\
B & Intermediate & 288 & -.0049 & .0252\\
B & Original practical & 222 & .0562 & .0773\\
B & Search-censored & 32 & .1957 & .1958\\
\bottomrule
\end{tabular}
\end{table}
Below identification, all 208 Layer-A and 112 Layer-B comparisons are identical programs; above finite value thresholds the corresponding counts are 384 and 192.
In the original practical subset, 323/446 Layer-A and 176/222 Layer-B cells have strictly positive improvement lower bounds, while 101 and 33 have strictly negative upper bounds; the remaining 22 and 13 are inconclusive.
These include regime-specific losses even when the stratum-average improvement is positive.
No new admission rule was executed and no old gate was replaced.
The old probe-length routine used the sufficient exponent $\log(0.5/0.025)$, giving an equal-prior average-error target of $0.025$ rather than a directional guarantee of $0.025$.
The reanalysis uses each probe's exact affinity $\exp(-I_B)$ as a valid bound on either directional error; over all 104 frozen probe designs these bounds range from $0.0352271$ to $0.0497465$.
Failure of this sufficient directional certificate does not establish that the actual directional errors exceed $0.025$.
No outcomes or probe lengths are changed to obtain a smaller error bound.
This outcome-known reanalysis gives population enclosures, not new confidence intervals for empirical means.

\section{Complete Threshold-Sandwich Scale Calculation}
For the independently initialized Gaussian running mean, let $c_a=\max\{(h+q)/\beta_m,b/\beta_e\}$ and $r=\lambda^b$.
The exact horizon is $N_a=s/c_a+O(1)$, with bounded terminal charges and rounding.
Main Eq.~\eqref{eq:mean-risk} and Eq.~\eqref{eq:geometric-correlation-sum} give $K_j(a)=c_a\{q^{-1}+\theta_j(1+r)/(1-r)\}$ and a uniform $O(s^{-2})$ remainder.
Removing probe charge $v$ changes $N_a$ by $O(v+1)$ and hence risk by $O((v+1)/s^2)$ for $s\ge s_0(v+1)$; compact persistence and the finite catalogue make these constants uniform.
Thus the phase-theorem risk conditions have a concrete Gaussian realization, separately from the constant-step affine certificate.
Write $d_\delta=\operatorname{kl}(1-\delta,\delta)$, $r_0=1-\gamma<1$, and $Q=\max{\cal Q}$.
Under either true instance, predictable spacing gives $\|C_{S_k}\|_{L^2}\le r_0\|C_{S_{k-1}}\|_{L^2}+\sqrt{\theta_+}$, hence a uniform bound $M_C=\sqrt{\theta_+}/(1-r_0)$.
The normalized observation $Y_k/\sqrt{q_k}$ has $L^2$ norm at most $M_C+1$.
For either candidate filter, write its posterior mean as $(1-t_{j,k})m^-_{j,k}+t_{j,k}Y_k/\sqrt{q_k}$, with predictable $0\le t_{j,k}\le1$.
Prediction multiplies this mean by a factor of magnitude at most $r_0$; the triangle inequality bounds both predictive means uniformly in $L^2$ by $M_F=r_0(M_C+1)/(1-r_0)$, starting from zero.
Both predictive variances are at most $\theta_+$, so $1\le V_{j,k}\le V_+=1+Q\theta_+$.
Main Eq.~\eqref{eq:conditional-kl} gives the expected, rather than pointwise, bound $\E_j\mathcal I_{j,1-j}(Z_k,A_k)\le i_+$, where $i_+=(\log V_++V_++4QM_F^2)/2$ is finite.
Extend a stopped rule predictably after stopping; nonnegativity and its deterministic budget bound $n_{\max}(s)\le c_Bs$ imply $\E_j\sum_{k\le\tau}\mathcal I_{j,1-j}\le i_+n_{\max}(s)$.
Together with main Eq.~\eqref{eq:info-constraint}, this proves $B_N\ge c_Nd_\delta$.
The fixed probe has information at least $i_->0$, so $n_\delta=\lceil i_-^{-1}\log(1/\delta)\rceil$ observations suffice and cost at most $c_pn_\delta+D_{\max}$.
Continuity of $\log(1/\delta)/d_\delta$ gives $n_\delta\le c_0+c_1d_\delta\le c_0+c_2B_N$.
Write $u_\delta\le c_pn_\delta+D_{\max}$ for the full probe charge.
The assumed expansion and budget-shift inequality give a high-instance improvement over the full-budget fallback certificate of at least $g_{\min}/s-\{2M+L(u_\delta+1)\}/s^2-\delta w_+/s$.
It is strictly positive for $s>\{4M+2L(u_\delta+1)\}/g_{\min}$, also requiring $s\ge s_0(u_\delta+1)$.
On instance zero, the assumed coefficient ordering gives $\epsilon_{\rm safe}\le\{2M+L(u_\delta+1)\}/s^2+\delta_0w_+/s$.
Thus $s\ge\sqrt{2\{2M+L(u_\delta+1)\}/\bar\epsilon}+2\delta_0w_+/\bar\epsilon$ suffices.
Using $2\sqrt{xy}\le x+y$, all three requirements are bounded by a class-constant multiple of $n_\delta+B_{\rm oracle}+B_\epsilon+1$.
Taking the maximum proves the upper bound in main Eq.~\eqref{eq:threshold-sandwich}; its lower bound follows because every admitted $\delta$-correct probe obeys both information constraints.

\section{Robust-Certificate Sensitivity Details}
This section supplies the proof of main Corollary~\ref{cor:robust-certificate}.
Write $\delta^+=C_q(r_q^+)^b$.
On the coverage event, $\delta^+\ge\delta$ and monotonicity on $\mu-2L\delta^+>0$ gives $a^+\ge a$ and $\beta^+\ge\beta$.
Evaluating the delayed Young inequality at the shared $\lambda^+=\sqrt{h/a^+}$ proves the upper recursion; separately optimized forcing terms need not be monotone.
For $\Delta\delta=\delta^+-\delta$, $\Delta K=K_q^+-K_q$, and $\Delta\Omega=\Omega_q^+-\Omega_q$,
\begin{align}
\Delta a
&=2\eta^2\Delta K+(2\eta L+4\eta^2L^2)\Delta\delta,\nonumber\\
\Delta\beta
&=2\eta^2\Delta\Omega+4\eta^2G^2\Delta\delta\nonumber\\
&\quad+4\eta G^2\left[
\frac{(\delta^+)^2}{\mu-2L\delta^+}
-\frac{\delta^2}{\mu-2L\delta}
\right],\nonumber\\
\Delta c
&=(\sqrt{a^+}+\sqrt h)^2-(\sqrt a+\sqrt h)^2.
\end{align}
For $h>0$,
\begin{align}
\Delta d
&\le(1+\sqrt{h/a^+})\Delta\beta
+\sqrt h\left(\frac1{\sqrt a}-\frac1{\sqrt{a^+}}\right)\beta\nonumber\\
&\quad+\frac{\sqrt{a^+}-\sqrt a}{\sqrt h}g,
\end{align}
and $\Delta d=\Delta\beta$ when $h=0$.
For $U(c,d)=d/(1-c)+c^n(R_0-d/(1-c))_+$, $|\partial_dU|\le(1-\bar c)^{-1}$ and $|\partial_cU|\le nR_0+\bar d(1-\bar c)^{-2}$.
The mean-value bound proves main Eq.~\eqref{eq:robust-score-gap}; applying it to the upper- and exact-score minimizers gives $U_{\widehat a}\le U_{a^\star}+2\max_a\varepsilon_a$.

\section{Stochastic-Delay Interpretation}
The controlled stochastic delay differential equation associated with one fixed resource action is
\begin{equation}
dE(t)=-A_qE(t-\tau(t))dt
+\sqrt{\eta}\,\Sigma_q^{1/2}(E_t)dW(t).
\end{equation}
Here $\Sigma_q(E_t)$ retains both additive innovation covariance and the state-dependent quadratic variation of random temporal-difference Jacobians.
For $V(E)=\|E\|^2$, It\^o's formula gives
\begin{equation}
{\cal L}_aV(E_t)
=-2E(t)^\top A_qE(t-\tau(t))
+\eta\tr\Sigma_q(E_t).
\end{equation}
Participation controls spatial covariance, spacing controls mixing, and $\eta$ controls their scale; the discrete theorem retains finite steps, integer delays and rounding.
\end{document}
